\documentclass[11pt]{article}
\usepackage[margin=2.5cm]{geometry}
\usepackage{booktabs}
\usepackage{array}
\usepackage{float}
\usepackage{url}

\title{Progress report\\ \large Plate-specific recognizer}
\date{20 October 2026}

\begin{document}
\maketitle

\noindent Code: \url{https://github.com/S1norin/inno-detecting-license-plates}

\bigskip
\textbf{Contributors}

\begin{center}
\begin{tabular}{@{}ll@{}}
\toprule
Name & Email \\
\midrule
Mikhail Krylov & m.krylov@innopolis.university \\
Maximilian Mifsud Bonici & m.mifsudbonici@innopolis.university \\
Andrei Polevoi & a.polevoi@innopolis.university \\
Evgenii Khovalyg & e.khovalyg@innopolis.university \\
\bottomrule
\end{tabular}
\end{center}

\section{New approach}

We trained CRNN recognizer for license plates. Convolutional layers read the
plate image, two bidirectional LSTM layers track the sequence of characters,
and a CTC output turns the sequence into text. It replaces EasyOCR from
week 1.

\begin{itemize}
\item What it is: a CRNN + BiLSTM + CTC text recognizer, about 3 million
  parameters
\item What it replaces: the EasyOCR recognition stage
\item The code is in \texttt{models/crnn/}
\end{itemize}

\section{Motivation}

In week 1 the detector was strong, but the recognizer was weak. EasyOCR read
only 35\% of characters correctly and got 0 out of 800 whole plates right on
the test set. The problem was that a general-purpose OCR model simply does
not transfer well to plate crops.

So we trained a model specifically on plates. We expected character accuracy
to rise well above 0.35 and exact-plate accuracy to rise from 0.

\section{Implementation}

\begin{table}[H]
\centering
\begin{tabular}{@{}lll@{}}
\toprule
Aspect & First approach & Second approach \\
\midrule
Model & EasyOCR 1.7.2, pretrained & CRNN: 4 conv blocks, 2 BiLSTM layers, CTC head \\
Training data & none (used as-is) & 3{,}500 original-scale train crops from CCPD \\
Training & none & 50 epochs, Adam lr 0.001, batch 32, seed 42, CPU only, $\sim$38 min \\
Character set & built-in, general text & the 58 characters from the training plates \\
Input & engine-specific & one grayscale crop, 48 $\times$ 256, aspect ratio kept \\
Evaluation & same harness, two protocols & unchanged \\
\bottomrule
\end{tabular}
\caption{Main differences from the first approach.}
\end{table}

\section{Initial results}

Test split, both engines through the same harness:

\begin{table}[H]
\centering
\begin{tabular}{@{}lrr@{}}
\toprule
Protocol / setting & Baseline (EasyOCR) & Second approach (CRNN) \\
\midrule
ocr\_only --- char.\ accuracy & 0.3474 & 0.7538 \\
ocr\_only --- exact plates & 0 / 800 & 269 / 800 \\
end\_to\_end --- char.\ accuracy & 0.3240 & 0.6561 \\
end\_to\_end --- exact plates & 0 / 800 & 133 / 800 \\
\bottomrule
\end{tabular}
\caption{Test results, second approach vs baseline.}
\end{table}

\begin{itemize}
\item The expected improvement happened: character accuracy more than
  doubled (34.7\% $\rightarrow$ 75.4\% on clean crops, 32.4\%
  $\rightarrow$ 65.6\% on the full pipeline), and exact plates went from 0
  to 269.
\item The full pipeline is still below the clean-crop number (65.6\% vs
  75.4\%): the detector misses 17 test plates, and a detector box is a
  rougher crop than the ground-truth box.
\item Large plates are read much worse than small ones (57\% vs 86\%),
  because we squash them into a 48~px tall input; the 10--19 degree tilt bin
  is the weakest (51\%, only 63 samples).
\end{itemize}

\section*{Remaining work}

\begin{itemize}
\item Fine-tune YOLO11 on our training split and compare it to the
  pretrained detector, then re-run end-to-end recognition with the better
  detector.
\item Try a taller input for large plates to improve their accuracy.
\item Final submission (10 November): failure-case analysis, runtime and
  memory numbers, reproducibility check on a clean clone, two-page summary,
  and demo.
\end{itemize}

\section*{References}

\begin{itemize}
\item CRNN details and reproduction: \texttt{models/crnn/REPORT.md}
\item Recognition baseline: \texttt{baseline/ocr/README.md}
\item Detection baseline: \texttt{baseline/detection/README.md}
\item Repository: \url{https://github.com/S1norin/inno-detecting-license-plates}
\end{itemize}

\end{document}
